\ifdefined\gkver\else\def\gkver{student}\fi
\documentclass[\gkver]{gaokaozhenti}
\begin{document}

\chapter{2006年四川理综}
%% paperType: 理综物理部分

\begin{enumerate}
\item
%% number: 14
%% paperName: 2006年四川理综第 14 题 6 分
%% typeId: 1
%% type: 单选题
%% chapter: 直线运动
%% point: 匀变速直线运动
%% method: 无
%% score: 6
%% degree: 850
%% duplicateId: 0
%% body:
2006 年我国自行研制的“枭龙”战机 04 架在四川某地试飞成功。假设该战机起飞前从静止开始做匀加速直线运动，达到起飞速度 $v$ 所需时间为 $t$，则起飞前的运动距离为\xzanswer{B}

\fourchoices[answer=B]
{$vt$}
{$\frac{vt}{2}$}
{$2vt$}
{不能确定}

%% answer: B
%% memo:
\memoanswer{
战机运动的平均速度为 $\bar{v} = \frac{v}{2}$，得 $s = \bar{v}t = \frac{v}{2}t$，故 B 正确。
}

\item
%% number: 15
%% paperName: 2006年四川理综第 15 题 6 分
%% typeId: 1
%% type: 单选题
%% chapter: 光学
%% point: 光电效应
%% method: 无
%% score: 6
%% degree: 650
%% duplicateId: 0
%% body:
现有 a、b、c 三束单色光，其波长关系为 $\lambda_{a} > \lambda_{b} > \lambda_{c}$。用 b 光束照射某种金属时，恰能发生光电效应。若分别用 a 光束和 c 光束照射该金属，则可以断定\xzanswer{A}

\fourchoices[answer=A]
{a 光束照射时，不能发生光电效应}
{c 光束照射时，不能发生光电效应}
{a 光束照射时，释放出的光电子数目最多}
{c 光束照射时，释放出的光电子的最大初动能最小}

%% answer: A
%% memo:
\memoanswer{
由题意知波长 $\lambda_{a} > \lambda_{b} > \lambda_{c}$，则有频率 $\nu_{a} < \nu_{b}$（极限频率）$< \nu_{c}$，故 A 正确，B、C 错误；又由光电效应方程知 D 错误。
}

\item
%% number: 16
%% paperName: 2006年四川理综第 16 题 6 分
%% typeId: 1
%% type: 单选题
%% chapter: 原子和原子核
%% point: 核聚变
%% method: 无
%% score: 6
%% degree: 650
%% duplicateId: 0
%% body:
某核反应方程 \ce{^{2}_{1}H + ^{3}_{1}H -> ^{4}_{2}He + X}。已知 \ce{^{2}_{1}H} 的质量为 $2.0136\Uu$，\ce{^{3}_{1}H} 的质量为 $3.0180\Uu$，\ce{^{4}_{2}He} 的质量为 $4.0026\Uu$，X 的质量为 $1.0087\Uu$。则下列说法中正确的是\xzanswer{B}

\fourchoices[answer=B]
{X 是质子，该反应释放能量}
{X 是中子，该反应释放能量}
{X 是质子，该反应吸收能量}
{X 是中子，该反应吸收能量}

%% answer: B
%% memo:
\memoanswer{
由质量数和电荷数守恒得 X 是中子，反应后质量亏损必释放能量，故 B 正确。
}

\item
%% number: 17
%% paperName: 2006年四川理综第 17 题 6 分
%% typeId: 1
%% type: 单选题
%% chapter: 电磁感应
%% point: 切割磁感线电动势
%% method: 无
%% score: 6
%% degree: 450
%% duplicateId: 0
%% body:
如图所示，接有灯泡 L 的平行金属导轨水平放置在匀强磁场中，一导体杆与两导轨良好接触并做往复运动，其运动情况与弹簧振子做简谐运动的情况相同。图中 O 位置对应于弹簧振子的平衡位置，P、Q 两位置对应于弹簧振子的最大位移处。若两导轨的电阻不计，则\xzanswer{D}

\onepicture[w=6.0cm]{figs/17.png}

\fourchoices[answer=D]
{杆由 O 到 P 的过程中，电路中电流变大}
{杆由 P 到 Q 的过程中，电路中电流一直变大}
{杆通过 O 处时，电路中电流方向将发生改变}
{杆通过 O 处时，电路中电流最大}

%% answer: D
%% memo:
\memoanswer{
杆切割磁感线运动的速度大，感应电动势大，电路中电流大。由 O 到 P 的过程中速度减小，故 A 错误；杆通过 O 处时速度最大，故 B 错误，D 正确；杆的运动方向在 P、Q 处改变时，电流方向发生改变，故 C 错误。
}

\item
%% number: 18
%% paperName: 2006年四川理综第 18 题 6 分
%% typeId: 1
%% type: 单选题
%% chapter: 交流电
%% point: 变压器
%% method: 无
%% score: 6
%% degree: 650
%% duplicateId: 0
%% body:
如图所示，理想变压器原、副线圈匝数之比为 $20:1$，原线圈接正弦交流电源，副线圈接入“$220\UV$，$60\UW$”灯泡一只，且灯泡正常发光。则\xzanswer{C}

\onepicture[w=5.5cm]{figs/18.png}

\fourchoices[answer=C]
{电流表的示数为 $\frac{3\sqrt{2}}{220}\UA$}
{电源输出功率为 $1200\UW$}
{电流表的示数为 $\frac{3}{220}\UA$}
{原线圈端电压为 $11\UV$}

%% answer: C
%% memo:
\memoanswer{
由灯泡正常发光知，副线圈电压为 $220\UV$，又原、副线圈电压与匝数成正比，故 D 错误；输入、输出功率相等，故 B 错误；电流表的示数为有效值，故 A 错误，C 正确。
}

\item
%% number: 19
%% paperName: 2006年四川理综第 19 题 6 分
%% typeId: 2
%% type: 多选题
%% chapter: 气体性质
%% point: 物体的内能
%% method: 无
%% score: 6
%% degree: 650
%% duplicateId: 0
%% body:
对一定质量的气体，下列说法中正确的是\xzanswer{BD}

\fourchoices[answer=BD]
{温度升高，压强一定增大}
{温度升高，分子热运动的平均动能一定增大}
{压强增大，体积一定减小}
{吸收热量，可能使分子热运动加剧、气体体积增大}

%% answer: BD
%% memo:
\memoanswer{
由 $\frac{pV}{T}$ 为恒量得 A、C 错误；温度升高，分子热运动的平均动能增大，故 B 正确；吸收热量，气体体积增大对外做功，温度可能升高，热运动加剧，故 D 正确。
}

\item
%% number: 20
%% paperName: 2006年四川理综第 20 题 6 分
%% typeId: 2
%% type: 多选题
%% chapter: 电场
%% point: 电势能电势差电场力做功
%% method: 无
%% score: 6
%% degree: 650
%% duplicateId: 0
%% body:
带电粒子 M 只在电场力作用下由 P 点运动到 Q 点，在此过程中克服电场力做了 $2.6\times10^{-6}\UJ$ 的功。那么\xzanswer{AD}

\fourchoices[answer=AD]
{M 在 P 点的电势能一定小于它在 Q 点的电势能}
{P 点的场强一定小于 Q 点的场强}
{P 点的电势一定高于 Q 点的电势}
{M 在 P 点的动能一定大于它在 Q 点的动能}

%% answer: AD
%% memo:
\memoanswer{
因克服电场力做功，电势能增加，动能减小，故 A、D 正确；P、Q 两点的场强大小不能确定，故 B 错误；粒子电性未知，所以 P、Q 两点的电势高低不能判定，故 C 错误。
}

\item
%% number: 21
%% paperName: 2006年四川理综第 21 题 6 分
%% typeId: 1
%% type: 单选题
%% chapter: 运动和力
%% point: 牛顿第二定律
%% method: 无
%% score: 6
%% degree: 450
%% duplicateId: 0
%% body:
质量不计的弹簧下端固定一小球。现手持弹簧上端使小球随手在竖直方向上以同样大小的加速度 $a$（$a < g$）分别向上、向下做匀加速直线运动。若忽略空气阻力，弹簧的伸长分别为 $x_{1}$、$x_{2}$；若空气阻力不能忽略且大小恒定，弹簧的伸长分别为 $x_{1}'$、$x_{2}'$。则\xzanswer{C}

\fourchoices[answer=C]
{$x_{1}' + x_{1} = x_{2} + x_{2}'$}
{$x_{1}' + x_{1} < x_{2} + x_{2}'$}
{$x_{1}' + x_{2}' = x_{1} + x_{2}$}
{$x_{1}' + x_{2}' < x_{1} + x_{2}$}

%% answer: C
%% memo:
\memoanswer{
忽略空气阻力时，向上有 $kx_{1}-mg=ma$，向下有 $mg-kx_{2}=ma$，得 $x_{1}+x_{2}=\frac{2mg}{k}$；考虑空气阻力时，向上有 $kx_{1}'-mg-f=ma$，向下有 $mg-kx_{2}'-f=ma$，得 $x_{1}'+x_{2}'=\frac{2mg}{k}$，故 C 正确。
}

\item
%% number: 22
%% paperName: 2006年四川理综第 22 题 6 分
%% typeId: 4
%% type: 实验题
%% chapter: 直线运动
%% point: 游标卡尺和螺旋测微仪
%% method: 无
%% score: 6
%% degree: 650
%% duplicateId: 0
%% body:
\begin{enumerate}
	\item
	在“用单摆测定重力加速度”的实验中：

	①测摆长时，若正确测出悬线长 $l$ 和摆球直径 $d$，则摆长为 \tkanswer[2.0cm]{$l+\frac{d}{2}$}；

	②测周期时，当摆球经过 \tkanswer[1.5cm]{平衡} 位置时开始计时并计数 1 次，测出经过该位置 $N$ 次（约 $60\sim100$ 次）的时间为 $t$，则周期为 \tkanswer[2.0cm]{$\frac{2t}{N-1}$}。

	此外，请你从下列器材中选用所需器材，再设计一个实验，粗略测出重力加速度 $g$，并参照示例填写下表（示例的方法不能再用）：

	A．天平；B．刻度尺；C．弹簧秤；D．电磁打点计时器；E．带夹子的重锤；F．纸带；G．导线若干；H．铁架台；I．低压交流电源；J．低压直流电源；G．小车；K．螺旋测微器；M．斜面（高度可调，粗糙程度均匀）。

	\begin{center}
	\begin{tblr}{colspec={|c|c|X[c]|}, width=\linewidth, hlines}
	 & 所选器材（只填器材序号） & 简述实验方法（不要求写出具体步骤） \\
	示例 & B、D、E、F、G、H、I & 安装仪器，接通电源，让纸带随重锤竖直下落。用刻度尺测出所需数据，处理数据，得出结果 \\
	实验设计 & & \\
	\end{tblr}
	\end{center}

	\item
	在“测定金属的电阻率”实验中，需要测量金属丝的长度和直径。现用最小分度为 $1\Umm$ 的米尺测量金属丝长度，图中箭头所指位置是拉直的金属丝两端在米尺上相对应的位置，测得的金属丝长度为 \tkanswer[1.8cm]{972.0} $\Umm$。在测量金属丝直径时，如果受条件限制，身边只有米尺 1 把和圆柱形铅笔 1 支。如何较准确地测量金属丝的直径？请简述测量方法：\tkanswer[4.0cm]{}。

	\onepicture[w=10.0cm, align=right]{figs/22.png}
\end{enumerate}
%% answer:
\jdanswer{
\begin{enumerate}
	\item
	① $l+\frac{d}{2}$；② 平衡，$\frac{2t}{N-1}$；实验设计：A、C、E；用弹簧秤称出带夹子重锤的重力大小 $G$，再用天平测出其质量 $m$，则 $g=G/m$。或 B、D、F、G、I、K、M；安装仪器，接通电源，让纸带随小车一起沿斜面下滑，用刻度尺测出所需数据，改变斜面高度再测一次，利用两次数据，由牛顿第二定律算出结果。

	\item
	① $972.0\Umm$；② 在铅笔上紧密排绕金属丝 $N$ 匝，用米尺量出该 $N$ 匝金属丝的宽度 $D$，由此可以计算得出金属丝的平均直径为 $\frac{D}{N}$。
\end{enumerate}
}
%% memo:
\memoanswer{
（1）①摆长等于悬线长加上摆球半径，即 $l+\frac{d}{2}$；②摆球经过平衡位置时速度最大，在该位置开始计时误差最小；测出经过该位置 $N$ 次的时间 $t$，则完成 $\frac{N-1}{2}$ 次全振动，周期 $T=\frac{2t}{N-1}$。测量重力加速度的实验设计见答案（用弹簧秤和天平称量法，或用打点计时器和斜面测量法）。

（2）①由米尺读数得金属丝长度为 $972.0\Umm$；②用累积法，在铅笔上紧密排绕金属丝 $N$ 匝，用米尺量出该 $N$ 匝金属丝的宽度 $D$，则金属丝的平均直径 $d=\frac{D}{N}$。
}

\item
%% number: 23
%% paperName: 2006年四川理综第 23 题 16 分
%% typeId: 6
%% type: 计算题
%% chapter: 功和能
%% point: 机械能守恒定律
%% method: 无
%% score: 16
%% degree: 650
%% duplicateId: 0
%% body:
荡秋千是大家喜爱的一项体育活动。随着科技的迅速发展，将来的某一天，同学们也许会在其它星球上享受荡秋千的乐趣。假设你当时所在星球的质量为 $M$、半径为 $R$，可将人视为质点，秋千质量不计、摆长不变、摆角小于 $90^{\circ}$，万有引力常量为 $G$。那么，

\begin{enumerate}
	\item
	该星球表面附近的重力加速度 $g_{\text{星}}$ 等于多少？
	\item
	若经过最低位置的速度为 $v_{0}$，你能上升的最大高度是多少？
\end{enumerate}
%% answer:
\jdanswer{
\begin{enumerate}
	\item
	$g_{\text{星}} = \frac{GM}{R^{2}}$
	\item
	$h = \frac{R^{2}v_{0}^{2}}{2GM}$
\end{enumerate}
}
%% memo:
\memoanswer{
（1）设人的质量为 $m$，在星球表面附近的重力等于万有引力，有 $mg_{\text{星}}=\frac{GMm}{R^{2}}$，解得 $g_{\text{星}} = \frac{GM}{R^{2}}$。

（2）设人能上升的最大高度为 $h$，由功能关系得 $mg_{\text{星}}h=\frac{1}{2}mv_{0}^{2}$，解得 $h = \frac{R^{2}v_{0}^{2}}{2GM}$。
}

\item
%% number: 24
%% paperName: 2006年四川理综第 24 题 19 分
%% typeId: 6
%% type: 计算题
%% chapter: 电场
%% point: 带电粒子的直线运动
%% method: 无
%% score: 19
%% degree: 450
%% duplicateId: 0
%% body:
如图所示的电路中，两平行金属板 A、B 水平放置，两板间的距离 $d = 40\Ucm$。电源电动势 $E = 24\UV$，内电阻 $r = 1\UO$，电阻 $R = 15\UO$。闭合开关 S，待电路稳定后，将一带正电的小球从 B 板小孔以初速度 $v_{0} = 4\Ums$ 竖直向上射入板间。若小球带电量为 $q = 1\times10^{-2}\UC$，质量为 $m = 2\times10^{-2}\Ukg$，不考虑空气阻力。那么，

\begin{enumerate}
	\item
	滑动变阻器接入电路的阻值为多大时，小球恰能到达 A 板？
	\item
	此时，电源的输出功率是多大？
\end{enumerate}

\onepicture[w=6.0cm, align=right]{figs/24.png}
%% answer:
\jdanswer{
\begin{enumerate}
	\item
	$R_{\text{滑}} = 8\UO$
	\item
	$P_{\text{出}} = 23\UW$
\end{enumerate}
}
%% memo:
\memoanswer{
（1）小球进入板间后，受重力和电场力作用，且到 A 板时速度为零，设两板间电压为 $U_{AB}$，由动能定理得 $-mgd - qU_{AB} = 0 - \frac{1}{2}mv_{0}^{2}$，滑动变阻器两端电压为 $U_{\text{滑}}=U_{AB}=8\UV$。设通过滑动变阻器电流为 $I$，由欧姆定律得 $I = \frac{E - U_{\text{滑}}}{R + r} = 1\UA$，滑动变阻器接入电路的电阻为 $R_{\text{滑}} = \frac{U_{\text{滑}}}{I} = 8\UO$。

（2）电源的输出功率为 $P_{\text{出}} = I^{2}(R + R_{\text{滑}}) = 23\UW$。
}

\item
%% number: 25
%% paperName: 2006年四川理综第 25 题 20 分
%% typeId: 6
%% type: 计算题
%% chapter: 磁场
%% point: 洛伦兹力
%% method: 无
%% score: 20
%% degree: 450
%% duplicateId: 0
%% body:
如图所示，在足够大的空间范围内，同时存在着竖直向上的匀强电场和垂直纸面向里的水平匀强磁场，磁感应强度 $B = 1.57\UT$。小球 1 带正电，其电量与质量之比 $\frac{q_{1}}{m_{1}} = 4\UCkg$，所受重力与电场力的大小相等；小球 2 不带电，静止放置于固定的水平悬空支架上。小球 1 向右以 $v_{0} = 23.59\Ums$ 的水平速度与小球 2 正碰，碰后经过 $0.75\Us$ 再次相碰。设碰撞前后两小球带电情况不发生改变，且始终保持在同一竖直平面内。（取 $g = 10\Umsq$）问：

\begin{enumerate}
	\item
	电场强度 $E$ 的大小是多少？
	\item
	两小球的质量之比 $\frac{m_{2}}{m_{1}}$ 是多少？
\end{enumerate}

\onepicture[w=6.0cm, align=right]{figs/25.png}
%% answer:
\jdanswer{
\begin{enumerate}
	\item
	$E = 2.5\UNC$
	\item
	$\frac{m_{2}}{m_{1}} = 11$
\end{enumerate}
}
%% memo:
\memoanswer{
（1）由题意知，小球 1 所受的重力与电场力始终平衡，则有 $m_{1}g = q_{1}E$，解得 $E = 2.5\UNC$。

（2）相碰后小球 1 做匀速圆周运动，由牛顿第二定律得 $q_{1}v_{1}B = m_{1}\frac{v_{1}^{2}}{R_{1}}$，解得半径为 $R_{1}=\frac{m_{1}v_{1}}{q_{1}B}$，运动的周期为 $T=\frac{2\pi m_{1}}{q_{1}B}=1\Us$。两小球运动时间为 $t=0.75\Us=\frac{3}{4}T$，小球 1 只能逆时针经 $\frac{3}{4}$ 个圆周时与小球 2 再次相碰。第一次相碰后小球 2 做平抛运动，由平抛规律得 $h = R_{1} = \frac{1}{2}gt^{2}$，$L = R_{1} = v_{2}t$。两小球第一次碰撞前后动量守恒，以水平向右为正方向，由动量守恒得 $m_{1}v_{0} = -m_{1}v_{1} + m_{2}v_{2}$。由上述两式得 $v_{2} = 3.75\Ums$，由半径公式得 $v_{1} = \frac{q_{1}BR_{1}}{m_{1}} = 17.66\Ums$，故两小球质量之比 $\frac{m_{2}}{m_{1}} = \frac{v_{0} + v_{1}}{v_{2}} = 11$。
}

\end{enumerate}

\end{document}
